\section{Appendix E: Prompt Design}
\label{app:prompts}

\subsection{The prompt is the task statement}
There is no hand-written prompt wrapper in this benchmark. The prompt is the
task's \texttt{instruction.md}, delivered verbatim by the harness. Its structure
is fixed across tasks:

\begin{enumerate}[leftmargin=1.4em,itemsep=1pt]
\item a machine-readable canary comment (stripped before delivery);
\item \texttt{\# Original Contest Problem} --- the verbatim contest statement,
with each cited figure embedded inline as \texttt{![图 N](fig\_N.png)};
\item any results tables the statement requires, as plain text;
\item \texttt{\# 题面附图} / \texttt{\# Figures} --- a note listing the
in-container figure paths (e.g.\ \texttt{/root/data/figures/fig\_1.png\ldots});
\item \texttt{\# Required Benchmark Artifact} --- instructs the agent to solve
the problem, points at the mirrored attachments, names the required output path,
and points at the output schema;
\item the time limit line (\texttt{You have 14400 seconds\ldots}).
\end{enumerate}

\subsection{The artifact block (verbatim template)}
\begin{verbatim}
# Required Benchmark Artifact

Solve the contest problem above. The official attachments are mirrored at
`/root/data/attachments` (flat filenames; see
`/root/data/attachments/INDEX.md`) and also remain at their full corpus paths
under `/root/data/repo`. Use them whenever they are relevant. Write your final
structured answer to:

/root/results/<slug>_result.json

The result must be valid JSON and should follow the shape shown in
`/root/data/output_schema.json`. Include the numerical quantities, modeling
decisions, diagnostics, sensitivity or uncertainty results, conclusions for the
contest questions, and generated artifact paths needed to support your answer.
You may also write supplementary CSV, XLSX, or PNG files under `/root/results`,
but the JSON file above is the required artifact.

You have 14400 seconds to complete this task. Do not cheat by using online
solutions or hints specific to this task.
\end{verbatim}

\subsection{Harness wrapping}
Three harnesses deliver the same text differently:
\begin{itemize}[leftmargin=1.4em,itemsep=1pt]
\item \textbf{terminus-2} (harbor): prepends a fixed command-line system prompt
(``You are an AI assistant tasked with solving command-line tasks in a Linux
environment\ldots provide batches of shell commands\ldots format your response as
JSON with \texttt{commands:[{keystrokes,duration}]}'') and then appends
\texttt{Task Description:} followed by the statement.
\item \textbf{codex CLI}: the statement is the positional prompt
(\texttt{codex exec \ldots\ -{}- '<statement>'}).
\item \textbf{opencode CLI}: the statement is the positional prompt
(\texttt{opencode run \ldots\ -{}- '<statement>'}).
\end{itemize}

\subsection{Version caveat}
The \texttt{instruction.md} files were revised after some trials were recorded:
a figure-path section was added, and one MCM statement was re-sourced from the
official PDF. The prompt actually delivered in a given trial is the one captured
in that trial's \texttt{trial.log} / \texttt{trajectory.json}; for
reproducibility, cite those, not the current file. This is why we describe the
prompt structure above rather than quoting today's file as if it were what every
model saw.

\subsection{What the prompt does \emph{not} say}
Deliberately, the prompt gives no modeling hints, no method suggestions, and no
pointer to the scored fields beyond the schema. It also does not tell the agent
which sub-question is scored, except that the artifact is required and must
contain ``the numerical quantities\ldots conclusions for the contest questions''.
This keeps the benchmark a test of doing the whole modeling problem, at the cost
of the agent having to infer which numbers matter.
